\section{Skeletons: where the crossings must be}\label{sec:skeletons}

Theorem~X says how much information must reach macroscopic crossings. This section says where they can be. At a scale that is
a fixed fraction of $L$, a linear-time head can turn only boundedly often, so its coarse motion is one of finitely many words;
crossings force two turns with overlapping legs; and the expansion of Theorem~X can always be found on one leg of one tape. A leg is
monotone only at the coarse scale: inside it the head may move back and forth over distances up to the scale.

\subsection{Crossings are zigzags}

\begin{lemma}[Crossing revisits are zigzags]\label{lem:zigzag}
\formal{MacroscopicSkeleton.lean}{crossing_isHop_opposite_polarity, exists_crossingZigzag_of_crossingHopsGT2H, card_tapeCrossingLowers_ge_of_macroOverlap}
Let $y$ be a unit-speed walk with horizon $T$ and $H\ge1$.
\begin{enumerate}[label=(\alph*)]
\item Two crossing revisits $(v_1,u_1)$, $(v_2,u_2)$ with $v_1<v_2<u_1<u_2$ visit different cells and lie on opposite sides of them:
one stays strictly above its cell and the other strictly below.
\item If two crossing revisits both have excursion greater than $2H$, there are a sign $\varepsilon\in\{\pm1\}$ and steps
$s_0<s_1<s_2<s_3<T$ such that $h=\varepsilon y$ satisfies
\[
  h(s_1)-h(s_0)\ge2H+1,\quad h(s_1)-h(s_2)\ge2H+1,\quad h(s_3)-h(s_2)\ge2H+1,\quad
  \min(h(s_1),h(s_3))-\max(h(s_0),h(s_2))\ge1 .
\]
\item Conversely, for every $\varepsilon$ and $s_0<s_1<s_2<s_3<T$, the walk has at least $2N$ steps that are earlier ends of
crossing revisits of excursion greater than $2H$, where $N=\max\bigl(0,\min(h(s_1),h(s_3))-\max(h(s_0),h(s_2))-2H\bigr)$, whatever the
walk does between these steps.
\end{enumerate}
\end{lemma}

\begin{outline}
(a) A walk with unit steps cannot pass a cell without visiting it (\lean{isWalk_ivt}); a revisit therefore stays on one side of its
cell, and the side of the second revisit is forced by the fact that the first cell is visited inside it but its own cell is not.
(b) Take the extreme points of the two excursions. (c) Every level strictly inside the overlap is crossed upward between $s_0$ and
$s_1$, downward between $s_1$ and $s_2$ and upward again between $s_2$ and $s_3$; the last visits before the turns give, for each
level, one upper and one lower revisit of excursion greater than $2H$, and these cross.
\end{outline}

\subsection{Skeletons}

\begin{definition}[Skeleton at scale $\Lambda$]\label{def:skeleton}
\formal{MacroscopicSkeleton.lean}{milestones, tapeSkeleton, macroTurns, traceSkeletonTuple}
Let $y$ be a walk with $y(0)=0$ and let $\Lambda\ge1$. Scan the steps $t=0,1,\dots,T-1$ keeping a grid point $g\in\Z$, initially
$0$. Whenever $\abs{y(t)-g\Lambda}\ge\Lambda$, record a \emph{milestone} $(t,\sigma)$ with $\sigma=+$ if $y(t)>g\Lambda$ and
$\sigma=-$ otherwise, and move $g$ one step in the direction~$\sigma$. The \emph{skeleton} of $y$ at scale $\Lambda$ is the word
$\sigma_1\sigma_2\cdots$ of recorded signs (\lean{tapeSkeleton}). A \emph{leg} is a maximal run of equal signs, and the number of
\emph{turns} is the number of sign changes (\lean{macroTurns}). The skeleton of a trace is the $k$-tuple of the skeletons of its
work tapes.
\end{definition}

\begin{lemma}[Confinement and finiteness]\label{lem:finite}
\formal{MacroscopicSkeleton.lean}{epochGrid_confines_and_exact, tapeSkeleton_length_mul_le, card_boundedSkeletonTuple, traceSkeletonTuple_mem_boundedSkeletonTuple}
\begin{enumerate}[label=(\alph*)]
\item At every step $\abs{y(t)-g(t)\Lambda}<\Lambda$, where $g(t)$ is the current grid point, and $y(t)=g(t)\Lambda$ exactly at
each milestone. A skeleton of length $m$ has $m\Lambda\le T$.
\item Let $\tau$ be any trace with $k$ work tapes and $T\le C(L+1)$, with $d,b\ge1$, $C\ge1$, $\theta>0$ and $d\ge\lceil8/\theta\rceil$.
At scale $\Lambda=\lfloor\theta L/4\rfloor$ we have $\Lambda\ge1$, and every work tape's skeleton has at most
$\Mmax=\lfloor16C/\theta\rfloor$ letters. Hence the skeleton of $\tau$ belongs to the set of $k$-tuples of words of length at most
$\Mmax$ over $\{+,-\}$, which has exactly $(2^{\Mmax+1}-1)^k$ elements and does not depend on~$d$.
\end{enumerate}
\end{lemma}

\begin{outline}
(a) is an induction over the scan. (b) Each milestone after the first needs the head to move $\Lambda$ cells, so $m\Lambda\le T\le
2CL$, while $\Lambda\ge\theta L/8$ once $\theta L\ge8$; hence $m\le16C/\theta$. The count of words of length at most $\Mmax$ is
$\sum_{i\le\Mmax}2^i$.
\end{outline}

\begin{definition}[Crossing-capable skeletons]\label{def:capable}
\formal{MacroscopicSkeleton.lean}{IsScaleFreeCrossingCapable, CrossingCapableSkeletonTuple}
Let $s$ be a skeleton and let $\gamma(i)$ be its grid coordinate after $i$ letters. The skeleton is \emph{crossing-capable} if for
some sign $\varepsilon$ and indices $i_0<i_1<i_2<i_3\le\abs s$ the function $h=\varepsilon\gamma$ satisfies
\[
  h(i_1)-h(i_0)\ge3,\quad h(i_1)-h(i_2)\ge3,\quad h(i_3)-h(i_2)\ge3,\quad \min(h(i_1),h(i_3))-\max(h(i_0),h(i_2))\ge-1
\]
(\lean{IsScaleFreeCrossingCapable}): up three grid steps, down three, up three, with overlapping ranges. Such a skeleton has at
least two turns.
\end{definition}

\begin{proposition}[Which skeletons carry crossings]\label{prop:capable}
\formal{MacroscopicSkeleton.lean}{isCrossingCapableSkeleton_of_tapeCrossingLowers_nonempty, isScaleFreeCrossingCapable_of_two_mul_le, two_le_macroTurns_of_isScaleFreeCrossingCapable, tapeCrossingLowers_nonempty_of_isDetourProofSkeleton}
Let $y$ be a walk with $y(0)=0$, $\Lambda\ge1$ and $H\ge2\Lambda$. If the tape of $y$ has an earlier end of a crossing pair of revisits
of excursion greater than $2H$, then its skeleton at scale $\Lambda$ is crossing-capable; in particular it has at least nine
letters and at least two turns. Conversely, if for some sign and indices $i_0<i_1<i_2<i_3$ the skeleton's grid satisfies
$(\min(h(i_1),h(i_3))-\max(h(i_0),h(i_2)))\Lambda\ge2H+1$, then the tape has such crossing steps.
\end{proposition}

\begin{outline}
The first part combines \cref{lem:zigzag}(b) with the confinement of \cref{lem:finite}(a): the four extreme steps lie within
$\Lambda$ of grid points, which turns the inequalities with margin $2H+1\ge4\Lambda+1$ into the grid inequalities above
(\lean{isCrossingCapableSkeleton_of_tapeCrossingLowers_nonempty}, which proves a finer statement at every scale). The converse is
\cref{lem:zigzag}(c) applied at the milestone steps, where the walk is exactly on the grid.
\end{outline}

Between the two conditions there is a narrow band of skeletons that may or may not carry crossings, depending on the motion inside
the legs; this band is non-empty, so no condition on the skeleton alone decides whether crossings occur.

\begin{lemma}[No macroscopic revisit within a leg]\label{lem:legs}
\formal{MacroscopicSkeleton.lean}{legOfEpoch_lt_of_isHop_gt2H}
Let $y$ be a walk with $y(0)=0$, $\Lambda\ge1$ and $H\ge2\Lambda$. Every revisit of excursion greater than $2H$ starts in an earlier
leg of the skeleton than the one in which it ends.
\end{lemma}

\begin{outline}
Within one leg the grid point moves monotonically. A revisit returns to its cell, so if it started and ended in the same leg it
would be confined to at most two adjacent epochs, with excursion less than $3\Lambda\le2H$.
\end{outline}

\subsection{Theorem S}

For an oblivious linear-time transposer $M$ with time constant $C$ and $\alpha>0$, put $\theta=\thetaM(M,C,\alpha)$,
$\Lambda^*=\lfloor\theta L/4\rfloor$, $H^*=\lfloor\theta L\rfloor$ and $\Mmax=\lfloor16C/\theta\rfloor$, and for a work tape $j$ and
$r\ge0$ let $\CL^{j,r}$ be the set of steps of $\CL_{H^*}$ that are earlier ends of revisits on tape $j$ and lie in the $r$-th leg of
tape $j$'s skeleton at scale $\Lambda^*$ (\lean{legCrossingLowers}).

\begin{theorem}[Theorem~S, skeletons and a single leg]\label{thm:S}
\formal{MacroscopicSkeleton.lean}{obliviousLinTime_reduces_to_crossingCapableSkeleton_and_singleLeg, obliviousLinTime_infinite_pinned_skeleton_and_leg}
Let $M$ be a machine with $k$ work tapes, $C\ge1$ and $0<\alpha\le1$, and let $d\ge\max(d_0,\lceil8/\theta\rceil)$ with $d_0$ as in
\cref{thm:X}. Suppose $\mathrm{ObliviousLinTime}(M,d,C,T)$.
\begin{enumerate}[label=(\alph*)]
\item $1\le\Lambda^*$ and $2\Lambda^*\le H^*$, and the skeleton of the common trace at scale $\Lambda^*$ belongs to the explicit
finite family of $k$-tuples of words of length at most $\Mmax$ in which some tape's word is crossing-capable.
\item For every set $V_0$ of at least $\alpha L$ input positions there are a work tape $j$, whose skeleton is crossing-capable
with at least nine letters and two turns, and a leg $r\le\Mmax$ such that every set separating $\Out(V_0)$ from $\CL^{j,r}$ has at
least
\[
  \frac{\cE(M,C,\alpha/2)}{2k(\Mmax+1)}\,\abs{V_0}
\]
elements.
\item If $\mathrm{ObliviousLinTime}(M,d,C,T_d)$ holds for all $d$ in an infinite set $S$, then there are one tuple of words, one tape
and one leg such that, for infinitely many $d\in S$, the skeleton is that tuple and (b) holds with that tape and leg for $V_0$ the
set of all input positions (with $\alpha=1$).
\end{enumerate}
\end{theorem}

\begin{outline}
(a) is \cref{lem:finite}(b) together with \cref{prop:capable}, since $\CL_{H^*}$ is non-empty by \cref{cor:many}. (b) By
\cref{lem:legs} each step of $\CL_{H^*}$ lies in a slot $(j,r)$ with $r\le\Mmax$, and there are at most $k(\Mmax+1)$ slots. If every
slot had a separating set smaller than the stated bound, their union would separate $\Out(V_0)$ from all of $\CL_{H^*}$
(\lean{isMengerVertexCut_biUnion}) with fewer than $\frac12\cE\abs{V_0}$ elements, contradicting \cref{thm:X}. (c) The skeleton,
tape and leg take finitely many values, independent of $d$, so one of them recurs infinitely often
(\lean{exists_infinite_fiber_of_mapsTo_finset}).
\end{outline}

The tape and the leg in (b) may depend on $V_0$ and on~$d$; (c) removes the dependence on $d$ along a subsequence. The finite
family in (a) is astronomically large, because $\Mmax=\lfloor16C/\theta\rfloor$ inherits the size of~$\theta$; what matters is that
it does not grow with~$d$. The Lean proves only that the crossing-capable tuples form a proper subfamily of all tuples (the idle
tuple is excluded, \lean{card_crossingCapableSkeletonTuple_lt}), not how much smaller it is.

\begin{remark}[Numerical: what the skeleton does not see]\label{rem:skeleton-limits}
Inside a leg the head may move back and forth over distances up to $\Lambda^*$, which is a fixed fraction of $L$. The skeleton does
not see this motion, and the explicit cascades of \cref{sec:cascades} hide all but boundedly many of their refinement stages inside
it. Our measurements on those schedules, and on a one-pass shear that \cref{thm:J} excludes, found skeletons independent of $d$ with a
single leg carrying linear expansion ($d\le128$). So, as far as these measurements go, the skeleton condition does not separate these
schedules from a transposer.
\end{remark}
